\documentclass[11pt]{article}
\usepackage[margin=1in]{geometry}
\usepackage{amsmath,amssymb,mathtools}
\usepackage{fontspec,unicode-math}
\setmainfont{Latin Modern Roman}
\setsansfont{Latin Modern Sans}
\setmathfont{Latin Modern Math}
\usepackage{xurl,hyperref}
\hypersetup{colorlinks=true,urlcolor=blue,linkcolor=blue}
\newcommand{\sref}[2]{\hyperlink{source:#1:#2}{[S#2]}}
\newcommand{\cataloguescope}[1]{\paragraph{Recorded mathematical question.}#1}
\begin{document}
\section{Converse implications among SIRSN properties}\label{statement}

\subsection{Definitions and mathematical statement}

\paragraph{Shared notation.}$\mathbb N=\{1,2,\ldots\}$, and $\mathbb Z,\mathbb Q,\mathbb R,\mathbb C$ denote the integers, rationals, reals and complex numbers. Any other conventions are specified locally.


A scale-invariant random spatial network (SIRSN) is a jointly specified family of random routes $R(z,z')$ between points of $\mathbb R^2$. Routes are non-self-intersecting curves of finite length made of straight segments, with finitely or countably many segments whose only possible limiting endpoints are $z,z'$. The reverse route is $R(z',z)$. Any two routes meeting twice coincide between those meeting points. For every finite set of endpoints, the spanning subnetwork is the union of all its pairwise routes. These finite-dimensional laws are consistent under restriction, measurable in the endpoints, and invariant under translations, rotations and simultaneous multiplication of all coordinates by any $c>0$.
Let $\Xi(\lambda)$ be independent planar Poisson points of intensity $\lambda$, coupled increasingly in $\lambda$ by Poisson rain, and let $S(\lambda)$ be the union of the routes between their pairs. These subnetworks have almost surely finite length in every bounded set. For a stationary edge set $F$, its edge intensity $\operatorname{int}(F)$ is defined by $\mathbb E\operatorname{len}(F\cap A)=\operatorname{int}(F)\operatorname{area}(A)$ for bounded Borel sets $A$; length counts each road segment only once. Write $B(z,r)$ for a Euclidean disc, $e=(1,0)$, and
\[
D_1=\operatorname{len}R(0,e),\quad \Delta=\mathbb E D_1,\quad \ell=\operatorname{int}S(1),
\]
\[
E(\lambda,r)=\bigcup_{\xi,\xi'\in\Xi(\lambda)}
\bigl(R(\xi,\xi')\setminus(B(\xi,r)\cup B(\xi',r))\bigr),\quad
E(\infty,r)=\bigcup_{\lambda>0}E(\lambda,r),\quad
p(r)=\lim_{\lambda\to\infty}\operatorname{int}E(\lambda,r).
\]
The SIRSN assumptions include $1<\Delta<\infty$ and $p(1)<\infty$; they imply $\ell<\infty$. A weak SIRSN replaces the requirement $p(1)<\infty$ by $\ell<\infty$. Routes are given primitives and need not be shortest paths for Euclidean length.
In a sampled network enlarged by a fixed starting point $z$, a singly infinite geodesic from $z$ is an unbounded path whose finite initial portions agree with initial portions of routes from $z$ to sampled points. Let $G_z(\lambda)$ count distinct such paths and $G_z=\lim_{\lambda\to\infty}G_z(\lambda)$. The unique-geodesic property $U$ is $G_z=1$ almost surely, for each fixed $z$. Under $U$, let $C$ be the assertion that the singly infinite geodesics from $0$ and $e$ eventually coincide.
For $0<\varepsilon<1<B$, let $Q(\lambda,\varepsilon,B)$ be the probability that the routes from all sampled points in $B(0,\varepsilon)$ to all sampled points outside $B(0,B)$ do not all first exit $B(0,1)$ at one common point. Put $Q(\infty,\varepsilon,B)=\lim_{\lambda\to\infty}Q(\lambda,\varepsilon,B)$, and define
\[F:\quad \lim_{\varepsilon\downarrow0,\,B\uparrow\infty}Q(\infty,\varepsilon,B)=0.\]
For $0<\eta<\delta<1/2$, truncate each route between sampled points in $B(0,\eta)$ and $B(e,\eta)$ at its first exit from $B(0,\delta)$ and its last entrance into $B(e,\delta)$. Let $\Psi(\lambda,\eta,\delta)$ be the probability that these middle subroutes do not all coincide, and put $\Psi(\infty,\eta,\delta)=\lim_{\lambda\to\infty}\Psi(\lambda,\eta,\delta)$. Define
\[H:\quad \lim_{\eta\downarrow0}\Psi(\infty,\eta,\delta)=0\quad\text{for every }0<\delta<1/2.\]
The source numbers $\ell<\infty$, $p(1)<\infty$, $C$, $F$, $H$ respectively as (16), (20), (49), (50), (51).
\paragraph{Statement recorded in the supplied source.}
Prove each implication or construct a counterexample within the corresponding network class: (i) a weak SIRSN must satisfy $p(1)<\infty$; (ii) $U$ implies $C$ for SIRSNs; (iii) $C$ implies $F$ for SIRSNs with $U$; (iv) $H$ implies $F$ for SIRSNs. In (iii), coalescence is considered with the unique-geodesic assumption under which the source defines it. \sref{statement}{2}
\subsection{Short English statement}
Decide all four proposed converses concerning finite road intensity, uniqueness and coalescence of infinite routes, and continuity of nearby routes.
\subsection{Sources}
\begin{description}
\item[\hypertarget{source:statement:1}{S1}] ULAM.AI and the UnsolvedMath Contributors, \emph{UnsolvedMath: A Curated Collection of Open Mathematics Problems}, catalogue entry. CC BY 4.0. \url{https://huggingface.co/datasets/ulamai/UnsolvedMath}.
\item[\hypertarget{source:statement:2}{S2}] David J. Aldous, \emph{Scale-Invariant Random Spatial Networks}, arXiv:1204.0817v1 (2012), Sections 6.5--7.2 and Section 8.4.1, Open Problem 32, pp. 41--46 and 49. \url{https://arxiv.org/pdf/1204.0817}.
\end{description}
\label{end:statement}
\end{document}
